\subsection{Estructura local de ciertos funtores cohomológicos}
\label{subsection:III.6.10}
\label{III.6.10}

\begin{proposition}[6.10.1]
\label{III.6.10.1}
Sea $S=\Spec(A)$ un esquema afín, y sea $Y^{(i)}$
$(1\leq i\leq n)$ una familia finita de $S$-esquemas afines,
planos sobre $S$. Para cada $i$, sea
\[
  f_i:X^{(i)}\longrightarrow Y^{(i)}
\]
un $S$-morfismo separado y cuasicompacto, y sea
$\sh P_\bullet^{(i)}$ un complejo acotado inferiormente de
$\sh O_{X^{(i)}}$-módulos cuasicoherentes. Sea $Y$ el producto
sobre $S$ de los esquemas $Y^{(i)}$.

Existe un complejo $\sh R_\bullet$ de $\sh O_Y$-módulos
cuasicoherentes, planos sobre $S$, con la propiedad siguiente.
Para todo $S$-esquema afín $S'$ y todo complejo acotado
inferiormente $\sh Q_\bullet'$ de $\sh O_{S'}$-módulos
cuasicoherentes, existe un isomorfismo
\begin{env}[6.10.1.1]
\label{III.6.10.1.1}
\[
  \mathbf{Tor}_\bullet^S
  \bigl(
    f_1,\ldots,f_n,1_{S'};
    \sh P_\bullet^{(1)},\ldots,\sh P_\bullet^{(n)},
    \sh Q_\bullet'
  \bigr)
  \xrightarrow{\sim}
  \sh H_\bullet
  \bigl(
    \sh R_\bullet\otimes_{\sh O_S}\sh Q_\bullet'
  \bigr),
  \tag{6.10.1.1}\label{III.6.10.1.1.eq}
\]
\end{env}
que es un isomorfismo de $\partial$-funtores en
$\sh Q_\bullet'$.

Además, para todo $S$-morfismo $u:S''\to S'$ de $S$-esquemas
afines, el diagrama
\begin{env}[6.10.1.2]
\label{III.6.10.1.2}
\[
\xymatrix@C=12pt@R=30pt{
  {\begin{gathered}
    \mathbf{Tor}_\bullet^S
    (f_1,\ldots,f_n,1_{S'};\\
    \sh P_\bullet^{(1)},\ldots,\sh P_\bullet^{(n)},
    \sh Q_\bullet')
  \end{gathered}}
  \ar[r]^-{\sim}\ar[d]
  &
  {\sh H_\bullet
    (\sh R_\bullet\otimes_{\sh O_S}\sh Q_\bullet')}
  \ar[d]
  \\
  {\begin{gathered}
    \mathbf{Tor}_\bullet^S
    (f_1,\ldots,f_n,1_{S''};\\
    \sh P_\bullet^{(1)},\ldots,\sh P_\bullet^{(n)},
    u^*(\sh Q_\bullet'))
  \end{gathered}}
  \ar[r]^-{\sim}
  &
  {\sh H_\bullet
    (\sh R_\bullet\otimes_{\sh O_S}u^*(\sh Q_\bullet'))}
}
\tag{6.10.1.2}\label{III.6.10.1.2.eq}
\]
\end{env}
es conmutativo. Aquí las flechas verticales son los
$(1_Y\times_Su)$-morfismos canónicos definidos en
\sref{III.6.7.10}.
\end{proposition}

\begin{proof}
\oldpage[III]{172}
Calculemos el hipertor mediante el método de \sref{III.6.6.2},
teniendo en cuenta la Observación \sref{III.6.6.6}. Con las
notaciones allí utilizadas, cada $L_{\bullet\bullet}^{(i)}$ es un
bicomplejo de $A_i$-módulos, donde $A_i$ denota el anillo de
$Y^{(i)}$. Tiene, por tanto, una resolución proyectiva de
Cartan--Eilenberg $M_{\bullet\bullet\bullet}^{(i)}$ (en el sentido
de \sref[0]{0.11.7.1}) por $A_i$-módulos proyectivos. Por
\sref{III.6.6.6}, el miembro izquierdo de
\eqref{III.6.10.1.1.eq} es canónicamente isomorfo a
\[
  \HH_\bullet
  \bigl(
    M_{\bullet\bullet\bullet}\otimes_AQ_\bullet'
  \bigr),
\]
donde
\[
  M_{\bullet\bullet\bullet}
  =
  M_{\bullet\bullet\bullet}^{(1)}
  \otimes_A
  M_{\bullet\bullet\bullet}^{(2)}
  \otimes_A\cdots\otimes_A
  M_{\bullet\bullet\bullet}^{(n)}
\]
y $\sh Q_\bullet'=\widetilde{Q_\bullet'}$.

Por hipótesis, los anillos $A_i$ son $A$-módulos planos. Por
tanto, los $M_{\bullet\bullet\bullet}^{(i)}$ son tricomplejos de
$A$-módulos planos \sref[0]{0.6.2.1}, y lo mismo ocurre con
$M_{\bullet\bullet\bullet}$. Además, si $B$ es el anillo de $Y$,
es decir, el producto tensorial sobre $A$ de los $A_i$, entonces
$M_{\bullet\bullet\bullet}$ es un tricomplejo de $B$-módulos.
En consecuencia, el complejo
\[
  \sh R_\bullet
  =
  \widetilde{M_{\bullet\bullet\bullet}},
\]
donde $M_{\bullet\bullet\bullet}$ se considera como complejo
simple, tiene las propiedades requeridas, como se deduce
inmediatamente de \sref{III.6.7.10}.
\end{proof}

\begin{corollary}[6.10.2]
\label{III.6.10.2}
En la Proposición \sref{III.6.10.1}, el complejo
$\sh R_\bullet$ puede tomarse acotado inferiormente. Si los
$\sh P_\bullet^{(i)}$ están acotados superiormente y los
$Y^{(i)}$ tienen dimensión cohomológica finita, entonces
$\sh R_\bullet$ puede tomarse acotado superiormente.
\end{corollary}

\begin{proof}
La primera afirmación se deduce de que los tres grados de cada
$M_{\bullet\bullet\bullet}^{(i)}$ están acotados inferiormente.
Por otra parte, si los anillos $A_i$ tienen dimensión
cohomológica finita, el tercer grado de cada
$M_{\bullet\bullet\bullet}^{(i)}$ solo toma un número finito de
valores, y por su construcción lo mismo ocurre con su primer grado
\sref{III.6.6.2}. Su segundo grado está acotado superiormente
exactamente cuando lo está el grado de $\sh P_\bullet^{(i)}$
\sref{III.6.6.2}, lo que demuestra la segunda afirmación.
\end{proof}

\begin{remarks}[6.10.3]
\label{III.6.10.3}
\textnormal{(i)}
Con las notaciones de \sref{III.6.10.1},
\[
  \sh H_\bullet
  (\sh R_\bullet\otimes_S\sh Q_\bullet')
  \simeq
  \mathbf{Tor}_\bullet(\sh R_\bullet,\sh Q_\bullet'),
\]
porque $\sh R_\bullet$ consta de $\sh O_Y$-módulos planos sobre
$S$ \sref{III.6.5.9}. Es, por tanto, según \sref{III.6.5.4}, el
objeto al que converge una sucesión espectral regular cuyo término
$E^2$ es
\begin{env}[6.10.3.1]
\label{III.6.10.3.1}
\[
  E^2_{pq}
  =
  \bigoplus_{q'+q''=q}
  \shTor_p^S
  \bigl(
    \sh H_{q'}(\sh R_\bullet),
    \sh H_{q''}(\sh Q_\bullet')
  \bigr).
  \tag{6.10.3.1}\label{III.6.10.3.1.eq}
\]
\end{env}
Esta es precisamente la sucesión espectral (e) de cambio de base
de \sref{III.6.9.3}.

\medskip
\textnormal{(ii)}
Sea $\sh R_\bullet'$ otro complejo de $\sh O_Y$-módulos
cuasicoherentes, planos sobre $S$, y sea
\[
  g:\sh R_\bullet'\longrightarrow\sh R_\bullet
\]
un homomorfismo de complejos para el cual
\[
  \sh H_\bullet(g):
  \sh H_\bullet(\sh R_\bullet')
  \longrightarrow
  \sh H_\bullet(\sh R_\bullet)
\]
es un isomorfismo. Por \sref{III.6.3.3} y
\sref{III.6.5.9}, $g$ induce un isomorfismo de
$\partial$-funtores en $\sh Q_\bullet'$,
\[
  \sh H_\bullet
  (\sh R_\bullet'\otimes_S\sh Q_\bullet')
  \xrightarrow{\sim}
  \sh H_\bullet
  (\sh R_\bullet\otimes_S\sh Q_\bullet'),
\]
tal que el diagrama
\[
\xymatrix@C=28pt@R=30pt{
  \sh H_\bullet
  (\sh R_\bullet'\otimes_S\sh Q_\bullet')
  \ar[r]^-{\sim}\ar[d]
  &
  \sh H_\bullet
  (\sh R_\bullet\otimes_S\sh Q_\bullet')
  \ar[d]
  \\
  \sh H_\bullet
  (\sh R_\bullet'\otimes_Su^*(\sh Q_\bullet'))
  \ar[r]^-{\sim}
  &
  \sh H_\bullet
  (\sh R_\bullet\otimes_Su^*(\sh Q_\bullet'))
}
\]
es conmutativo. Así, las propiedades de \sref{III.6.10.1} no
determinan por completo el complejo $\sh R_\bullet$.

\medskip
\oldpage[III]{173}
\textnormal{(iii)}
En la demostración de \sref{III.6.10.1}, los
$M_{\bullet\bullet\bullet}^{(i)}$ pueden elegirse formados por
$A_i$-módulos libres, como se deduce inmediatamente dualizando la
demostración de \sref[0]{0.11.5.2.1}. Entonces los
\[
  M_{\bullet\bullet\bullet}^{(i)}\otimes_{A_i}B
\]
constan de $B$-módulos libres; como
$M_{\bullet\bullet\bullet}$ es su producto tensorial sobre $B$,
se puede exigir además en \sref{III.6.10.1} que
$\sh R_\bullet$ esté asociado a un complejo de $B$-módulos
libres.

Además, por (M, XVII, 1.2), el tricomplejo
$M_{\bullet\bullet\bullet}$ depende funtorialmente de cada uno de
los bicomplejos $L_{\bullet\bullet}^{(i)}$ y, por tanto, de cada
$\sh P_\bullet^{(i)}$, una vez fijado un recubrimiento finito de
cada $X^{(i)}$. Aquí se entienden por «morfismos» de tricomplejos
las clases de homomorfismos módulo homotopía. Sustituir un
recubrimiento de $X^{(i)}$ por otro más fino proporciona, para
$L_{\bullet\bullet}^{(i)}$, homomorfismos definidos salvo
homotopía por \sref{III.6.6.8}; así, con la convención precedente
sobre los morfismos, $M_{\bullet\bullet\bullet}$ es un funtor en
cada $\sh P_\bullet^{(i)}$.

Precisaremos esta dependencia funtorial, en particular el
comportamiento de $\sh R_\bullet$ respecto de sucesiones exactas
cortas
\[
  0\longrightarrow\sh P_\bullet^{\prime(i)}
  \longrightarrow\sh P_\bullet^{(i)}
  \longrightarrow\sh P_\bullet^{\prime\prime(i)}
  \longrightarrow0,
\]
en el capítulo dedicado a un álgebra general de funtores
cohomológicos mencionado en \sref{III.6.1.3}.
\end{remarks}

\begin{scholium}[6.10.4]
\label{III.6.10.4}
El hecho de que $\sh R_\bullet$ conste de $\sh O_Y$-módulos
planos sobre $S$ implica inmediatamente que
\[
  \sh H_\bullet
  (\sh R_\bullet\otimes_S\sh Q_\bullet')
\]
es un funtor homológico en $\sh Q_\bullet'$; compárese el
razonamiento de \sref{III.7.7.1}. Como se señaló en
\sref{III.6.1.1}, esta propiedad motiva la introducción del
hipertor.

Pongamos
\[
  X=X^{(1)}\times_SX^{(2)}\times_S\cdots\times_SX^{(n)},
  \qquad
  f=f_1\times_Sf_2\times_S\cdots\times_Sf_n,
\]
\[
  \sh P_\bullet
  =
  \sh P_\bullet^{(1)}
  \otimes_S\sh P_\bullet^{(2)}
  \otimes_S\cdots\otimes_S\sh P_\bullet^{(n)},
\]
y
\[
  X'=X\times_SS',
  \qquad
  Y'=Y\times_SS',
  \qquad
  f'=f\times_S1_{S'}.
\]
Los problemas de cambio de base llevan a estudiar la
hipercohomología
\[
  \sh H^*_{f'}
  (\sh P_\bullet\otimes_S\sh N')
\]
como funtor del $\sh O_{S'}$-módulo cuasicoherente $\sh N'$, o,
equivalentemente, la hipercohomología
\[
  \sh H^*_f
  (\sh P_\bullet\otimes_S\sh N)
\]
como funtor del $\sh O_S$-módulo cuasicoherente $\sh N$. Cuando
los $\sh P_\bullet^{(i)}$, y por tanto $\sh P_\bullet$, son
planos sobre $S$, los resultados precedentes y
\sref{III.6.7.6} muestran que este es efectivamente un funtor
cohomológico en $\sh N$. Esto deja de ser cierto cuando no se
supone la planitud de los $\sh P_\bullet^{(i)}$, y el estudio de
\[
  \sh H^*_f(\sh P_\bullet\otimes_S\sh N)
\]
ya no puede abordarse mediante los métodos usuales del álgebra
homológica.

Utilizaremos principalmente el caso $n=1$, $Y=S$, en el que
$\sh P_\bullet$ consta de $\sh O_X$-módulos planos sobre $Y$.
En este caso se tiene el teorema siguiente.
\end{scholium}

\begin{theorem}[6.10.5]
\label{III.6.10.5}
Sea $Y=\Spec(A)$ un esquema afín noetheriano, sea
$f:X\to Y$ un morfismo propio y sea $\sh P_\bullet$ un complejo
acotado inferiormente de $\sh O_X$-módulos coherentes, planos
sobre $Y$. Existe un complejo acotado inferiormente
$\sh L_\bullet$ de $\sh O_Y$-módulos, cuyos términos son de la
forma
\[
  \sh L_i=\sh O_Y^{m_i},
\]
y un isomorfismo
\begin{env}[6.10.5.1]
\label{III.6.10.5.1}
\[
  \sh H^*
  (f,\sh P_\bullet\otimes_Y\sh Q_\bullet)
  \xrightarrow{\sim}
  \sh H_\bullet
  (\sh L_\bullet\otimes_Y\sh Q_\bullet)
  \tag{6.10.5.1}\label{III.6.10.5.1.eq}
\]
\end{env}
de $\partial$-funtores en el complejo acotado inferiormente
$\sh Q_\bullet$ de $\sh O_Y$-módulos cuasicoherentes.

Además, para todo morfismo $u:Y'\to Y$, pongamos
\[
  X'=X_{(Y')},
  \qquad
  f'=f_{(Y')},
  \qquad
  \sh P_\bullet'
  =
  \sh P_\bullet\otimes_Y\sh O_{Y'},
  \qquad
  \sh L_\bullet'=u^*(\sh L_\bullet).
\]
\oldpage[III]{174}
Entonces $\sh L_\bullet'$ es un complejo de
$\sh O_{Y'}$-módulos localmente libres de tipo finito, y existe
un isomorfismo
\begin{env}[6.10.5.2]
\label{III.6.10.5.2}
\[
  \sh H^*
  (f',\sh P_\bullet'\otimes_{Y'}\sh Q_\bullet')
  \xrightarrow{\sim}
  \sh H_\bullet
  (\sh L_\bullet'\otimes_{Y'}\sh Q_\bullet')
  \tag{6.10.5.2}\label{III.6.10.5.2.eq}
\]
\end{env}
de $\partial$-funtores en el complejo acotado inferiormente
$\sh Q_\bullet'$ de $\sh O_{Y'}$-módulos cuasicoherentes. Estos
isomorfismos pueden elegirse de modo que el diagrama
\begin{env}[6.10.5.3]
\label{III.6.10.5.3}
\[
\xymatrix@C=24pt@R=30pt{
  \sh H^*
  (f,\sh P_\bullet\otimes_Y\sh Q_\bullet)
  \ar[r]^-{\sim}\ar[d]
  &
  \sh H_\bullet
  (\sh L_\bullet\otimes_Y\sh Q_\bullet)
  \ar[d]
  \\
  \sh H^*
  (f',\sh P_\bullet'\otimes_{Y'}u^*(\sh Q_\bullet))
  \ar[r]^-{\sim}
  &
  \sh H_\bullet
  (\sh L_\bullet'\otimes_{Y'}u^*(\sh Q_\bullet))
}
\tag{6.10.5.3}\label{III.6.10.5.3.eq}
\]
\end{env}
sea conmutativo.
\end{theorem}

\begin{proof}
La Proposición \sref{III.6.10.1} da primero un complejo acotado
inferiormente $\sh R_\bullet$ \sref{III.6.10.2} de
$\sh O_Y$-módulos cuasicoherentes libres sobre $Y$
\sref{III.6.10.3}~(iii), junto con un isomorfismo
\begin{env}[6.10.5.4]
\label{III.6.10.5.4}
\[
  \sh H^*
  (f,\sh P_\bullet\otimes_Y\sh Q_\bullet)
  \xrightarrow{\sim}
  \sh H_\bullet
  (\sh R_\bullet\otimes_Y\sh Q_\bullet)
  \tag{6.10.5.4}\label{III.6.10.5.4.eq}
\]
\end{env}
de $\partial$-funtores en $\sh Q_\bullet$. Los términos de
$\sh R_\bullet$ no son necesariamente, \emph{a priori},
$\sh O_Y$-módulos de tipo finito.

Apliquemos \eqref{III.6.10.5.4.eq} cuando $\sh Q_\bullet$ es el
complejo concentrado en un grado cuyo único término es $\sh O_Y$.
Se sigue que $\sh H_\bullet(\sh R_\bullet)$ es isomorfo a
$\sh H^*(f,\sh P_\bullet)$ y, por tanto, consta de
$\sh O_Y$-módulos coherentes \sref{III.6.2.5}. Por
\sref[0]{0.11.9.2}, existe entonces un complejo acotado
inferiormente $\sh L_\bullet$ cuyos términos están asociados a
$A$-módulos libres de tipo finito, y un homomorfismo
\[
  \sh L_\bullet\longrightarrow\sh R_\bullet
\]
que induce un isomorfismo
\[
  \sh H_\bullet(\sh L_\bullet)
  \xrightarrow{\sim}
  \sh H_\bullet(\sh R_\bullet).
\]
La fórmula \eqref{III.6.10.5.1.eq} se deduce ahora de
\sref{III.6.10.3}~(ii).

Las demás afirmaciones de \sref{III.6.10.5} se deducen de
\sref{III.6.10.1} y \sref{III.6.10.3}~(ii) cuando $Y'$ es afín.
En el caso general, recubrámoslo por abiertos afines $(V_\alpha)$.
Si \eqref{III.6.10.5.2.eq} se forma sobre cada $V_\alpha$, sus
restricciones desde $V_\alpha$ y $V_\beta$ a un abierto afín
\[
  W\subset V_\alpha\cap V_\beta
\]
coinciden ambas con el isomorfismo formado sobre $W$. Esto se
deduce de la conmutatividad de \eqref{III.6.10.1.2.eq}, aplicada
a las inclusiones canónicas $W\to V_\alpha$ y $W\to V_\beta$, y
demuestra que los isomorfismos locales se pegan.
\end{proof}

\begin{remark}[6.10.6]
\label{III.6.10.6}
En los capítulos siguientes aplicaremos \sref{III.6.10.5}
principalmente cuando $\sh P_\bullet$ está concentrado en un grado
y su único término es un $\sh O_X$-módulo coherente $\sh F$ plano
sobre $Y$. Como
\[
  \sh H_n(\sh L_\bullet)
  =
  \sh H^{-n}(f,\sh F)
  =
  R^{-n}f_*(\sh F)
  \qquad\text{\sref{III.6.2.1}},
\]
los $\sh H_n(\sh L_\bullet)$ se anulan para $n>0$. Veremos más
adelante, en \sref{III.7.7.12}~(i), que $\sh L_\bullet$ puede
elegirse con términos solo en grados $\leq0$ y, por tanto, con un
número finito de términos, siempre que la exigencia de que sus
términos estén asociados a $A$-módulos libres de tipo finito se
sustituya por la de que sean localmente libres de tipo finito.
El complejo $\sh L_\bullet$ correspondiente a tal
$\sh O_X$-módulo $\sh F$ no parece tener ninguna propiedad
especial
\oldpage[III]{175}
aparte de la restricción precedente sobre sus grados. Cabe, pues,
preguntarse recíprocamente si, dado un complejo acotado
inferiormente $\sh L_\bullet$ cuyos términos son
$\sh O_Y$-módulos asociados a $A$-módulos proyectivos de tipo
finito y cuyos términos de grado $>0$ se anulan, existen un
$Y$-esquema $X$, proyectivo y plano sobre $Y$, y un
$\sh O_X$-módulo localmente libre $\sh F$ tales que exista un
isomorfismo
\[
  \sh H^*
  (f,\sh F\otimes_Y\sh Q_\bullet)
  \xrightarrow{\sim}
  \sh H_\bullet
  (\sh L_\bullet\otimes_Y\sh Q_\bullet)
\]
funtorial en $\sh Q_\bullet$. Un resultado de este tipo reduciría
por completo la teoría cohomológica de los módulos coherentes
planos sobre $Y$ en $Y$-esquemas propios a la teoría, «salvo
homotopía», de los complejos de $A$-módulos proyectivos de tipo
finito sobre un anillo noetheriano $A$.
\end{remark}
